\documentclass[10pt,a4paper]{amsart}
\usepackage[margin=19mm]{geometry}
\usepackage{amsmath,amssymb,mathtools}
\usepackage[scr=rsfs,cal=euler]{mathalfa}
\usepackage{xfrac}
\usepackage[unicode,hidelinks,bookmarks=false]{hyperref}
\newcommand{\ie}{i.e.\ }
\newcommand{\cf}{cf.\ }
\newcommand{\resp}{respectively\ }
\newcommand{\Subsection}{\subsection}
\providecommand{\nobreakdash}{\nobreak\hbox{-}}
\setlength{\emergencystretch}{2em}
\multlinegap0pt
\usepackage{amssymb}
\newtheorem{prop}{Proposition}
\title{Asymptotic expansions of integrals and Nielsen's polylogarithms}
\author{Markus Kuba, Moti Levy}
\date{Source: arXiv:2604.05895v1 (CC BY 4.0)}
\begin{document}
\maketitle

\noindent\emph{Source and licence.} Asymptotic expansions of integrals and Nielsen's polylogarithms. Version arXiv:2604.05895v1 (CC BY 4.0), distributed by arXiv under the Creative Commons Attribution 4.0 International licence, http://creativecommons.org/licenses/by/4.0/, as recorded in the arXiv licence metadata for this exact version and shown on the abstract page https://arxiv.org/abs/2604.05895v1. Full licence notice: \texttt{licenses/arxiv-2604.05895-CC-BY-4.0.txt}.\par\smallskip

\noindent\textit{Editorial scope.} Counted unit: the article's prop Prop:asympMom (source0.tex lines 308--327). The article's complete original proof of it is delivered with it (source0.tex lines 358--417, one \texttt{proof} environment of 60 lines, which the article places away from the statement). No mathematics is written, repaired or re-derived: the statement and the proof are the article's own lines, in the article's own notation, and no retained line is substituted or re-flowed.  No further source-local context is needed: every cross-reference of every retained line resolves inside the two delivered spans.  Nothing was omitted from any retained span, so the package carries no omission record.  No retained line contains a citation, so the wrapper carries no bibliography and no bibliography entry.  No retained line contains a figure environment, an image-inclusion command or drawing code, so no image is delivered and the package declares no asset dependency.  Editorial formatting only, in addition: an AMS \texttt{amsart} preamble that restates the article's own declarations and macros used by the retained lines, namely the environments \texttt{prop} on separate counters, together with the standard packages those lines need.  The retained environments print in this document as Proposition 1 in order of appearance, whereas the article prints the same items with the numbering of its own full document.  The complete native source of the article as distributed in the arXiv e-print for this exact version is bundled unchanged as \texttt{sources/197987/source0.tex}.
\begin{prop}
\label{Prop:asympMom} The moments $\varphi_{r}=\int_{0}^{1}f(u)u^{r}$ satisfy
the asymptotic expansion
\[
\varphi_{r}\sim\sum_{\nu=0}^{\infty}\frac{\beta_{\nu}}{r^{\nu+1}},\qquad
\beta_{\nu}=(-1)^{\nu}\sum_{\ell=0}^{\nu}%
%TCIMACRO{\QATOPD{\{}{\}}{\nu+1}{\ell+1}}%
%BeginExpansion
\genfrac{\{}{\}}{0pt}{}{\nu+1}{\ell+1}%
%EndExpansion
f^{(\ell)}(1),
\]
where $%
%TCIMACRO{\QATOPD{\{}{\}}{\nu}{\ell}}%
%BeginExpansion
\genfrac{\{}{\}}{0pt}{}{\nu}{\ell}%
%EndExpansion
$ denote the Stirling numbers of the second kind, also called Stirling
partition numbers.
\end{prop}

\begin{proof}
[Proof of Proposition~\ref{Prop:asympMom}]We use the substitution $u=e^{-t}$
to obtain a Laplace-type integral:
\[
\varphi_{r}=\int_{0}^{\infty}g(t)e^{-rt}\,dt,\quad g(t):=e^{-t}f(e^{-t}).
\]
Watson's lemma yields
\[
\varphi_{r}\sim\sum_{\nu=0}^{\infty}\frac{g^{(\nu)}(0)}{r^{\nu+1}}.
\]
In order to find a simple expression for $g^{(\nu)}(0)$ we use an operator
formula. Let $\partial_{z}$ denote the differential operator with respect to
$z$, $\theta_{z}=z\partial_{z}$ the theta or homogeneity differential operator
and $E_{z=c}$ the evaluation operator at $z=c$. We observe that
\[
\partial_{t}e^{-t}f(e^{-t})=E_{x=e^{-t}}(-\theta_{x})xf(x),
\]
such that by induction
\[
g^{(\nu)}(t)=E_{x=e^{-t}}(-1)^{\nu}\theta_{x}^{\nu}xf(x).
\]
Next we use the expansion
\[
\theta_{x}^{\nu}=\sum_{k=0}^{\nu}%
%TCIMACRO{\QATOPD{\{}{\}}{\nu}{k}}%
%BeginExpansion
\genfrac{\{}{\}}{0pt}{}{\nu}{k}%
%EndExpansion
x^{k}\partial_{x}^{k}%
\]
to obtain
\[
g^{(\nu)}(t)=(-1)^{\nu}\sum_{k=0}^{\nu}%
%TCIMACRO{\QATOPD{\{}{\}}{\nu+1}{k+1}}%
%BeginExpansion
\genfrac{\{}{\}}{0pt}{}{\nu+1}{k+1}%
%EndExpansion
e^{-(k+1)t}f^{(k)}(e^{-t}),
\]
where we have used the basic recurrence relation
\[%
%TCIMACRO{\QATOPD{\{}{\}}{\nu+1}{k}}%
%BeginExpansion
\genfrac{\{}{\}}{0pt}{}{\nu+1}{k}%
%EndExpansion
=k%
%TCIMACRO{\QATOPD{\{}{\}}{\nu}{k}}%
%BeginExpansion
\genfrac{\{}{\}}{0pt}{}{\nu}{k}%
%EndExpansion
+%
%TCIMACRO{\QATOPD{\{}{\}}{\nu}{k-1}}%
%BeginExpansion
\genfrac{\{}{\}}{0pt}{}{\nu}{k-1}%
%EndExpansion
,\quad0<k<\nu,
\]
for the Stirling numbers of the second kind. Finally, we evaluate at $t=0$ to
obtain the stated identity.
\end{proof}

\end{document}
